\documentclass{amsart}
% For dealing with references we use the comment environment
\usepackage{verbatim}
\newenvironment{reference}{\comment}{\endcomment}
%\newenvironment{reference}{}{}
\newenvironment{slogan}{\comment}{\endcomment}
\newenvironment{history}{\comment}{\endcomment}

% For commutative diagrams we use Xy-pic
\usepackage[all]{xy}

% We use 2cell for 2-commutative diagrams.
\xyoption{2cell}
\UseAllTwocells

% We use multicol for the list of chapters between chapters
\usepackage{multicol}

% This is generally recommended for better output
\usepackage{lmodern}
\usepackage[T1]{fontenc}

% For cross-file-references

% Package for hypertext links:
\usepackage{hyperref}

% For any local file, say "hello.tex" you want to link to please

% Theorem environments.
%
\theoremstyle{plain}
\newtheorem{theorem}[subsection]{Theorem}
\newtheorem{proposition}[subsection]{Proposition}
\newtheorem{lemma}[subsection]{Lemma}

\theoremstyle{definition}
\newtheorem{definition}[subsection]{Definition}
\newtheorem{example}[subsection]{Example}
\newtheorem{exercise}[subsection]{Exercise}
\newtheorem{situation}[subsection]{Situation}

\theoremstyle{remark}
\newtheorem{remark}[subsection]{Remark}
\newtheorem{remarks}[subsection]{Remarks}

\numberwithin{equation}{subsection}

% Macros
%
\def\lim{\mathop{\mathrm{lim}}\nolimits}
\def\colim{\mathop{\mathrm{colim}}\nolimits}
\def\Spec{\mathop{\mathrm{Spec}}}
\def\Hom{\mathop{\mathrm{Hom}}\nolimits}
\def\Ext{\mathop{\mathrm{Ext}}\nolimits}
\def\SheafHom{\mathop{\mathcal{H}\!\mathit{om}}\nolimits}
\def\SheafExt{\mathop{\mathcal{E}\!\mathit{xt}}\nolimits}
\def\Sch{\mathit{Sch}}
\def\Mor{\mathop{\mathrm{Mor}}\nolimits}
\def\Ob{\mathop{\mathrm{Ob}}\nolimits}
\def\Sh{\mathop{\mathit{Sh}}\nolimits}
\def\NL{\mathop{N\!L}\nolimits}
\def\CH{\mathop{\mathrm{CH}}\nolimits}
\def\proetale{{pro\text{-}\acute{e}tale}}
\def\etale{{\acute{e}tale}}
\def\QCoh{\mathit{QCoh}}
\def\Ker{\mathop{\mathrm{Ker}}}
\def\Im{\mathop{\mathrm{Im}}}
\def\Coker{\mathop{\mathrm{Coker}}}
\def\Coim{\mathop{\mathrm{Coim}}}

% Boxtimes
%
\DeclareMathSymbol{\boxtimes}{\mathbin}{AMSa}{"02}

%
% Macros for moduli stacks/spaces
%
\def\QCohstack{\mathcal{QC}\!\mathit{oh}}
\def\Cohstack{\mathcal{C}\!\mathit{oh}}
\def\Spacesstack{\mathcal{S}\!\mathit{paces}}
\def\Quotfunctor{\mathrm{Quot}}
\def\Hilbfunctor{\mathrm{Hilb}}
\def\Curvesstack{\mathcal{C}\!\mathit{urves}}
\def\Polarizedstack{\mathcal{P}\!\mathit{olarized}}
\def\Complexesstack{\mathcal{C}\!\mathit{omplexes}}
% \Pic is the operator that assigns to X its picard group, usage \Pic(X)
% \Picardstack_{X/B} denotes the Picard stack of X over B
% \Picardfunctor_{X/B} denotes the Picard functor of X over B
\def\Pic{\mathop{\mathrm{Pic}}\nolimits}
\def\Picardstack{\mathcal{P}\!\mathit{ic}}
\def\Picardfunctor{\mathrm{Pic}}
\def\Deformationcategory{\mathcal{D}\!\mathit{ef}}

\setlength{\emergencystretch}{3em}
\begin{document}
\title[Stacks Project, Tag 07QR]{348 --- Henselization preserves the G-ring property}
\maketitle
\noindent Copyright (C) 2005--2025 Johan de Jong. Stacks Project authors. Source: \href{https://stacks.math.columbia.edu/tag/07QR}{Tag 07QR}.

\noindent Editorial difficulty note (not author text). Difficulty H2: The proof compares individual completed formal fibres with products of henselization fibres. It handles strict henselization separately through a formally smooth completion map and separable residue-field extensions, tracking the base field of geometric regularity throughout..

\noindent Editorial provenance note (not author text). History: excerpt from revision \texttt{a0444\allowbreak 6e57e\allowbreak c1fbc\allowbreak 252a8\allowbreak 71afc\allowbreak ec775\allowbreak 2fb28\allowbreak 07b14}, more-algebra.tex, lines 12429--12491. Reference presentation was adapted; original-author context and core support blocks were retained or added. The mathematical wording is unchanged. Licensed under GNU FDL 1.2 or later; no invariant sections or cover texts. See ../licenses/COPYING. Context blocks reproduce author definitions and supporting results; they are not additional counted problems.

\begin{definition}
\label{definition-G-ring}
A ring $R$ is called a {\it G-ring} if $R$ is Noetherian and for every
prime $\mathfrak p$ of $R$ the ring map
$R_\mathfrak p \to (R_\mathfrak p)^\wedge$ is regular.
\end{definition}

\begin{lemma}
\label{lemma-henselization-G-ring}
Let $R$ be a Noetherian local ring which is a G-ring.
Then the henselization $R^h$ and the strict henselization $R^{sh}$
are G-rings.
\end{lemma}

\begin{proof}
We will use the criterion of Lemma \href{https://stacks.math.columbia.edu/tag/07PT}{lemma check G ring maximal ideals (Tag 07PT)}.
Let $\mathfrak q \subset R^h$ be a prime and set
$\mathfrak p = R \cap \mathfrak q$. Set $\mathfrak q_1 = \mathfrak q$
and let $\mathfrak q_2, \ldots, \mathfrak q_t$
be the other primes of $R^h$ lying over $\mathfrak p$, so that
$R^h \otimes_R \kappa(\mathfrak p) =
\prod\nolimits_{i = 1, \ldots, t} \kappa(\mathfrak q_i)$, see
Lemma \href{https://stacks.math.columbia.edu/tag/07QQ}{lemma fibres henselization (Tag 07QQ)}.
Using that $(R^h)^\wedge = R^\wedge$
(Lemma \href{https://stacks.math.columbia.edu/tag/06LJ}{lemma henselization noetherian (Tag 06LJ)}) we see
$$
\prod\nolimits_{i = 1, \ldots, t}
(R^h)^\wedge \otimes_{R^h} \kappa(\mathfrak q_i) =
(R^h)^\wedge \otimes_{R^h} (R^h \otimes_R \kappa(\mathfrak p)) =
R^\wedge \otimes_R \kappa(\mathfrak p)
$$
Hence $(R^h)^\wedge \otimes_{R^h} \kappa(\mathfrak q_i)$
is geometrically regular over $\kappa(\mathfrak p)$ by assumption.
Since $\kappa(\mathfrak q_i)$ is separable algebraic over $\kappa(\mathfrak p)$
it follows from Algebra, Lemma
\href{https://stacks.math.columbia.edu/tag/07QH}{lemma geometrically regular over separable algebraic (Tag 07QH)} that
$(R^h)^\wedge \otimes_{R^h} \kappa(\mathfrak q_i)$ is
geometrically regular over $\kappa(\mathfrak q_i)$.

\medskip\noindent
Let $\mathfrak r \subset R^{sh}$ be a prime and set
$\mathfrak p = R \cap \mathfrak r$. Set $\mathfrak r_1 = \mathfrak r$
and let $\mathfrak r_2, \ldots, \mathfrak r_s$
be the other primes of $R^{sh}$ lying over $\mathfrak p$, so that
$R^{sh} \otimes_R \kappa(\mathfrak p) =
\prod\nolimits_{i = 1, \ldots, s} \kappa(\mathfrak r_i)$, see
Lemma \href{https://stacks.math.columbia.edu/tag/07QQ}{lemma fibres henselization (Tag 07QQ)}.
Then we see that
$$
\prod\nolimits_{i = 1, \ldots, s}
(R^{sh})^\wedge \otimes_{R^{sh}} \kappa(\mathfrak r_i) =
(R^{sh})^\wedge \otimes_{R^{sh}} (R^{sh} \otimes_R \kappa(\mathfrak p)) =
(R^{sh})^\wedge \otimes_R \kappa(\mathfrak p)
$$
Note that $R^\wedge \to (R^{sh})^\wedge$ is formally smooth
in the $\mathfrak m_{(R^{sh})^\wedge}$-adic topology, see
Lemma \href{https://stacks.math.columbia.edu/tag/06LJ}{lemma henselization noetherian (Tag 06LJ)}.
Hence $R^\wedge \to (R^{sh})^\wedge$ is regular by
Proposition \href{https://stacks.math.columbia.edu/tag/07PM}{proposition fs regular (Tag 07PM)}.
We conclude that $(R^{sh})^\wedge \otimes_{R^{sh}} \kappa(\mathfrak r_i)$
is regular over $\kappa(\mathfrak p)$ by
Lemma \href{https://stacks.math.columbia.edu/tag/07QI}{lemma regular composition (Tag 07QI)} as
$R^\wedge \otimes_R \kappa(\mathfrak p)$ is regular over $\kappa(\mathfrak p)$
by assumption. Since $\kappa(\mathfrak r_i)$ is separable algebraic over
$\kappa(\mathfrak p)$
it follows from Algebra, Lemma
\href{https://stacks.math.columbia.edu/tag/07QH}{lemma geometrically regular over separable algebraic (Tag 07QH)} that
$(R^{sh})^\wedge \otimes_{R^{sh}} \kappa(\mathfrak r_i)$ is
geometrically regular over $\kappa(\mathfrak r_i)$.
\end{proof}
\end{document}
